\documentclass[11pt]{ctexart}
\usepackage{a4}
\usepackage{graphicx}
\usepackage{geometry}
\usepackage{amsmath}
\usepackage{mathtools}
\geometry{a4paper,left=2cm,right=2cm,top=2.5cm,bottom=2.5cm}

\title{几何基础作业5}
\author{无名}
\date{\today}

\begin{document}
\maketitle
\section{证明正方形边长与对角线不可公度}

\begin{figure}[htb]
    \begin{center}
        \includegraphics[width=0.5\linewidth]{1.png}
        \caption{}
    \end{center}
\end{figure}

如图1，设边长$AB=a=1$，对角线$AC=b=\sqrt{2}$，显然$b=\sqrt{2}a$

我们让$b-a$，在BD上截取$a$，那么剩余部分就是$r_1=b-a$。

由相似，我们得到$\frac{EF}{BC}=\frac{EG}{BD}$，从而新的小正方形边长为$a_1=r_1=b-a$

那么新的小正方形对角线长度为$b_1=\sqrt{2}r_1=2-\sqrt{2}$

若a与b可以公度，即$\sqrt{2}=p/q$，那么就应该存在一个u使得所有的$a_i,b_i$可以被u同时量尽， 从而得到无穷递减正整数序列$\{a_i/u\}$ 

显然这种序列并不存在，故若a与b可以公度的假设不成立，也即为不可公度。
\newpage
\section{证明正三角形的一条边与这条边上面的高不可公度}

\begin{figure}[htb]
    \begin{center}
        \includegraphics[width=0.5\linewidth]{2.png}
        \caption{}
    \end{center}
\end{figure}

如图2，我们有正三角形$ABC$，边长$a$，高$h=\frac{\sqrt{3}}{2}a$

现在在高线$AD$上取$AF=\frac{a}{2}$，那么剩余部分就是$FD=h-\frac{a}{2}$

显然此时$AGH$也是正三角形，它和$ABC$相似，边长$a_1=\frac{a}{\sqrt{3}}$，高$h_1=\frac{h}{\sqrt{3}}$

重复这一过程，构造正三角形$AKL$边长$a_2=\frac{a_1}{\sqrt{3}}$，高$h_2=\frac{h_1}{\sqrt{3}}$

以此类推我们得到正三角形序列

\[a_k=\frac{a}{(\sqrt{3})^k} \hspace{1em} h_k=\frac{h}{(\sqrt{3})^k}\]

假设$a$与$h$可公度，那么就应该存在一个$u$使得所有的$a_i,h_i$可以被u同时量尽， 从而得到无穷递减正整数序列$\{a_i/u\}$和$\{h_i/u\}$

显然这种序列并不存在，故若$a$与$h$可以公度的假设不成立，也即为不可公度。
\end{document}